% ===== PRÉAMBULE CANONIQUE — à coller VERBATIM en tête de CHAQUE .tex =====
\documentclass[11pt,a4paper]{article}
\usepackage[utf8]{inputenc}\usepackage[T1]{fontenc}\usepackage{lmodern}
\usepackage[french]{babel}
\usepackage{amsmath,amssymb,mathtools}\usepackage{icomma}
\usepackage[version=4]{mhchem}          % \ce{...} : équations & formules
\usepackage{siunitx}\sisetup{locale=FR} % \qty{}{}, \si{}, \ang{}
\usepackage{chemfig}                    % \chemfig{...} : structures développées
\usepackage{xcolor}\usepackage{colortbl}
\usepackage{tikz}\usepackage{pgfplots}\pgfplotsset{compat=1.18}
\usepackage[most]{tcolorbox}\usepackage{enumitem}\usepackage{array,booktabs}
\usepackage[a4paper,margin=2.2cm]{geometry}
\usetikzlibrary{arrows.meta,calc,angles,quotes,positioning,patterns,backgrounds,shapes.geometric}
\definecolor{primary}{RGB}{222,49,99}\definecolor{primarydark}{RGB}{150,22,55}
\definecolor{defcol}{RGB}{0,114,178}\definecolor{methcol}{RGB}{52,92,148}
\definecolor{excol}{RGB}{0,135,90}\definecolor{attncol}{RGB}{200,40,55}
\tcbset{boxbase/.style={enhanced,breakable,boxrule=0.9pt,arc=2.5pt,
  left=3mm,right=3mm,top=2.3mm,bottom=2.3mm,fonttitle=\bfseries,coltitle=white}}
% --- boîtes TUTORIEL (noms EXACTS, ne pas renommer) ---
\newtcolorbox{cle}[1][]{boxbase,colback=defcol!6,colframe=defcol,
  title={\textsc{À retenir}\ifx\relax#1\relax\relax\else\textmd{ : \itshape #1}\fi}}
\newtcolorbox{methode}[1][]{boxbase,colback=methcol!7,colframe=methcol,sharp corners,
  title={\textsc{Méthode}\ifx\relax#1\relax\relax\else\textmd{ --- \itshape #1}\fi}}
\newtcolorbox{exemplebox}[1][]{boxbase,colback=excol!5,colframe=excol,
  title={\textsc{Exemple}\ifx\relax#1\relax\relax\else\textmd{ : \itshape #1}\fi}}
\newtcolorbox{piege}[1][]{boxbase,colback=attncol!6,colframe=attncol,
  title={\textsc{Piège}\ifx\relax#1\relax\relax\else\textmd{ : \itshape #1}\fi}}
\newtcolorbox{remarque}[1][]{boxbase,colback=black!4,colframe=black!45,
  title={\textsc{Remarque}\ifx\relax#1\relax\relax\else\textmd{ : \itshape #1}\fi}}
% --- boîtes EXERCICE / CORRIGÉ (noms EXACTS, auto-comptées) ---
\newtcolorbox[auto counter]{exo}[1][]{boxbase,colback=primary!4,colframe=primary,
  title={\textsc{Exercice}~\thetcbcounter\ifx\relax#1\relax\relax\else\textmd{ --- \itshape #1}\fi}}
\newtcolorbox[auto counter]{corrige}[1][]{boxbase,colback=excol!4,colframe=excol,
  title={\textsc{Corrigé}~\thetcbcounter\ifx\relax#1\relax\relax\else\textmd{ --- \itshape #1}\fi}}
\newcommand{\R}{\mathbb{R}}\newcommand{\N}{\mathbb{N}}\newcommand{\Z}{\mathbb{Z}}\newcommand{\ds}{\displaystyle}
% (un \usepackage{fancyhdr}+en-tête est possible ; il est ignoré côté web)
% =========================================================================
\begin{document}
% THEME_TITLE: Notation scientifique et arrondi
\sisetup{per-mode=symbol}

Deux gestes d'écriture indissociables en chimie : mettre un nombre sous la forme $a\times10^{n}$, et le
\emph{réduire} au bon nombre de chiffres significatifs (CS) par un arrondi propre.

\begin{cle}[notation scientifique]
Un nombre est en \textbf{notation scientifique} lorsqu'il s'écrit
\[ a\times10^{n}\qquad\text{avec}\quad 1\le |a|<10\ \text{et}\ n\in\mathbb{Z}. \]
La \textbf{mantisse} $a$ (un seul chiffre non nul avant la virgule) porte \emph{tous} les CS ; la
\textbf{puissance de dix} ne fixe que l'ordre de grandeur et n'apporte aucun CS.
\end{cle}

\begin{methode}[mettre un nombre en notation scientifique]
\begin{enumerate}[label=\arabic*.,leftmargin=1.8em,topsep=2pt,itemsep=1pt]
  \item placer la virgule \emph{juste après le premier chiffre non nul} ;
  \item compter de combien de rangs la virgule a bougé : c'est $|n|$ ;
  \item $n>0$ si le nombre était «~grand~» ($\ge10$), $n<0$ s'il était «~petit~» ($<1$) ;
  \item ne garder dans la mantisse que les CS (on supprime les zéros de cadrage).
\end{enumerate}
\end{methode}

\begin{exemplebox}[aller-retour]
$\num{20960}=2,096\times10^{4}$ \quad(virgule reculée de $4$ rangs)\qquad
$\num{0,02096}=2,096\times10^{-2}$ \quad(virgule avancée de $2$ rangs). Dans les deux cas, $4$ CS.
\end{exemplebox}

\begin{cle}[arrondir au rang voulu]
Pour arrondir, on regarde le \textbf{premier chiffre supprimé} (juste à droite du dernier chiffre gardé) :
\begin{itemize}[leftmargin=1.6em,topsep=2pt,itemsep=1pt]
  \item $0,1,2,3,4$ : on arrondit \textbf{par défaut} (le dernier chiffre gardé ne change pas) ;
  \item $5,6,7,8,9$ : on arrondit \textbf{par excès} (on ajoute $1$ au dernier chiffre gardé).
\end{itemize}
\end{cle}

\begin{center}
\begin{tikzpicture}[font=\small]
  \draw[->,>=Stealth,thick,black!70] (0,0) -- (10.4,0);
  \foreach \x/\d in {0/0,1.15/1,2.3/2,3.45/3,4.6/4,5.75/5,6.9/6,8.05/7,9.2/8,10.35/9}{
     \draw (\x,0.12) -- (\x,-0.12);
     \node[below,font=\footnotesize] at (\x,-0.14) {\d};}
  \draw[excol,line width=2pt] (0,0.32) -- (4.6,0.32);
  \node[excol,above,font=\footnotesize] at (2.3,0.36) {arrondi \textbf{par défaut} (0--4)};
  \draw[attncol,line width=2pt] (5.75,0.32) -- (10.35,0.32);
  \node[attncol,above,font=\footnotesize] at (8.05,0.36) {arrondi \textbf{par excès} (5--9)};
\end{tikzpicture}
\end{center}

\begin{exemplebox}[arrondir à un nombre de CS donné]
$3,148$ à $3$ CS $\to 3,15$ (le $8$ supprimé fait arrondir par excès). \quad
$2,996$ à $3$ CS $\to 3,00$ : la retenue se propage «~en cascade~», et l'on \emph{garde} les deux zéros
pour conserver $3$ CS.
\end{exemplebox}

\begin{piege}[deux réflexes de survie]
\begin{itemize}[leftmargin=1.6em,topsep=2pt,itemsep=2pt]
  \item \textbf{Garder les zéros de queue} après arrondi : $2,996\to 3,00$ (et non $3$), sinon on perd des CS.
  \item \textbf{N'arrondir qu'à la fin} d'un calcul : on garde des \emph{chiffres de garde} aux étapes
        intermédiaires. Arrondir trop tôt propage l'erreur.
\end{itemize}
\end{piege}

\begin{remarque}[notation ingénieur]
Variante où l'exposant $n$ est un \emph{multiple de $3$}, pour coller aux préfixes du Système
International (kilo $10^{3}$, milli $10^{-3}$, micro $10^{-6}$\dots). Exemple :
$2,096\times10^{-2}\,\si{\litre}=20,96\times10^{-3}\,\si{\litre}=\qty{20,96}{\milli\litre}$.
\end{remarque}

\end{document}
